
\documentclass{article}
\usepackage{colortbl}
\usepackage{makecell}
\usepackage{multirow}
\usepackage{supertabular}

\begin{document}

\newcounter{utterance}

\centering \large Interaction Transcript for game `matchit\_ascii', experiment `similar\_grid\_2', episode 4 with claude{-}opus{-}4{-}8{-}azure.
\vspace{24pt}

{ \footnotesize  \setcounter{utterance}{1}
\setlength{\tabcolsep}{0pt}
\begin{supertabular}{c@{$\;$}|p{.15\linewidth}@{}p{.15\linewidth}p{.15\linewidth}p{.15\linewidth}p{.15\linewidth}p{.15\linewidth}}
    \# & $\;$A & \multicolumn{4}{c}{Game Master} & $\;\:$B\\
    \hline

    \theutterance \stepcounter{utterance}  
    & & \multicolumn{4}{p{0.6\linewidth}}{
        \cellcolor[rgb]{0.9,0.9,0.9}{
            \makecell[{{p{\linewidth}}}]{
                \texttt{\tiny{[P1$\langle$GM]}}
                \texttt{Вы участвуете в совместной игре на угадывание.} \\
\texttt{Цель — выяснить, совпадают ли эта сетка и другая картинка, которую вижу только я.} \\
\texttt{Пожалуйста, сначала опишите вашу сетку. Затем я дам своё описание, и мы сможем задавать друг другу вопросы о сетках, чтобы понять, одинаковые ли они. Теперь начните ваше краткое описание сетки с "ОПИСАНИЕ:", затем добавьте описание. Ничего больше не добавляйте. Вот ваша сетка:} \\
\texttt{X X X X X} \\
\texttt{X ▢ ▢ ▢ X} \\
\texttt{X ▢ ▢ ▢ X} \\
\texttt{X ▢ ▢ ▢ X} \\
\texttt{X X X X X} \\
            }
        }
    }
    & \\ \\

    \theutterance \stepcounter{utterance}  
    & \multicolumn{4}{p{0.6\linewidth}}{
        \cellcolor[rgb]{0.9,0.9,0.9}{
            \makecell[{{p{\linewidth}}}]{
                \texttt{\tiny{[P1$\rangle$GM]}}
                \texttt{ОПИСАНИЕ: Сетка 5×5. Внешняя рамка полностью заполнена символами X (все клетки по краям). Внутренний квадрат 3×3 в центре состоит из пустых клеток.} \\
            }
        }
    }
    & & \\ \\

    \theutterance \stepcounter{utterance}  
    & & & \multicolumn{2}{p{0.3\linewidth}}{
        \cellcolor[rgb]{0.9,0.9,0.9}{
            \makecell[{{p{\linewidth}}}]{
                \texttt{\tiny{[GM$|$GM]}}
                \texttt{continue} \\
            }
        }
    }
    & & \\ \\

    \theutterance \stepcounter{utterance}  
    & & \multicolumn{4}{p{0.6\linewidth}}{
        \cellcolor[rgb]{0.9,0.9,0.9}{
            \makecell[{{p{\linewidth}}}]{
                \texttt{\tiny{[GM$\rangle$P2]}}
                \texttt{Вы участвуете в совместной игре на угадывание.} \\
\texttt{Цель — выяснить, совпадают ли эта сетка и другая картинка, которую вижу только я.} \\
\texttt{Пожалуйста, сначала опишите вашу сетку. Затем я дам своё описание, и мы сможем задавать друг другу вопросы о сетках, чтобы понять, одинаковые ли они. Теперь начните ваше краткое описание сетки с "ОПИСАНИЕ:", затем добавьте описание. Ничего больше не добавляйте. Вот ваша сетка:} \\
\texttt{X ▢ X X X} \\
\texttt{X ▢ ▢ ▢ X} \\
\texttt{X ▢ ▢ ▢ X} \\
\texttt{X ▢ ▢ ▢ X} \\
\texttt{X ▢ X X X} \\
            }
        }
    }
    & \\ \\

    \theutterance \stepcounter{utterance}  
    & & & \multicolumn{4}{p{0.6\linewidth}}{
        \cellcolor[rgb]{0.9,0.9,0.9}{
            \makecell[{{p{\linewidth}}}]{
                \texttt{\tiny{[GM$\langle$P2]}}
                \texttt{ОПИСАНИЕ: Сетка 5×5. Крайние левый и правый столбцы полностью заполнены X. В верхней строке: X, пусто, X, X, X. В нижней строке: X, пусто, X, X, X. Второй столбец полностью пустой (все ячейки пустые). Центральная область 3×3 (строки 2–4, столбцы 3–4) пустая, кроме крайних X по бокам.} \\
            }
        }
    }
     \\ \\

    \theutterance \stepcounter{utterance}  
    & & & \multicolumn{2}{p{0.3\linewidth}}{
        \cellcolor[rgb]{0.9,0.9,0.9}{
            \makecell[{{p{\linewidth}}}]{
                \texttt{\tiny{[GM$|$GM]}}
                \texttt{continue} \\
            }
        }
    }
    & & \\ \\

    \theutterance \stepcounter{utterance}  
    & & \multicolumn{4}{p{0.6\linewidth}}{
        \cellcolor[rgb]{0.9,0.9,0.9}{
            \makecell[{{p{\linewidth}}}]{
                \texttt{\tiny{[GM$\rangle$P2]}}
                \texttt{Это моёОПИСАНИЕ: Сетка 5×5. Внешняя рамка полностью заполнена символами X (все клетки по краям). Внутренний квадрат 3×3 в центре состоит из пустых клеток.} \\
\texttt{Теперь задайте вопрос, чтобы выяснить новые аспекты моей сетки, которые могут отличаться от вашей сетки. Начните с "ВОПРОС:" и не добавляйте ничего больше.} \\
            }
        }
    }
    & \\ \\

    \theutterance \stepcounter{utterance}  
    & & & \multicolumn{4}{p{0.6\linewidth}}{
        \cellcolor[rgb]{0.9,0.9,0.9}{
            \makecell[{{p{\linewidth}}}]{
                \texttt{\tiny{[GM$\langle$P2]}}
                \texttt{ВОПРОС: В вашей сетке во второй строке первая клетка — X, а вторая клетка (столбец 2) заполнена X или пустая?} \\
            }
        }
    }
     \\ \\

    \theutterance \stepcounter{utterance}  
    & & & \multicolumn{2}{p{0.3\linewidth}}{
        \cellcolor[rgb]{0.9,0.9,0.9}{
            \makecell[{{p{\linewidth}}}]{
                \texttt{\tiny{[GM$|$GM]}}
                \texttt{continue} \\
            }
        }
    }
    & & \\ \\

    \theutterance \stepcounter{utterance}  
    & & \multicolumn{4}{p{0.6\linewidth}}{
        \cellcolor[rgb]{0.9,0.9,0.9}{
            \makecell[{{p{\linewidth}}}]{
                \texttt{\tiny{[P1$\langle$GM]}}
                \texttt{Это моёОПИСАНИЕ: Сетка 5×5. Крайние левый и правый столбцы полностью заполнены X. В верхней строке: X, пусто, X, X, X. В нижней строке: X, пусто, X, X, X. Второй столбец полностью пустой (все ячейки пустые). Центральная область 3×3 (строки 2–4, столбцы 3–4) пустая, кроме крайних X по бокам.} \\
\texttt{ВОПРОС: В вашей сетке во второй строке первая клетка — X, а вторая клетка (столбец 2) заполнена X или пустая?Начните свой ответ с "ОТВЕТ:" и не добавляйте ничего другого.} \\
            }
        }
    }
    & \\ \\

    \theutterance \stepcounter{utterance}  
    & \multicolumn{4}{p{0.6\linewidth}}{
        \cellcolor[rgb]{0.9,0.9,0.9}{
            \makecell[{{p{\linewidth}}}]{
                \texttt{\tiny{[P1$\rangle$GM]}}
                \texttt{ОТВЕТ: Во второй строке моей сетки вторая клетка (столбец 2) пустая. Однако замечу различие: в моей верхней строке все клетки — X (X X X X X), тогда как у вас в верхней строке вторая клетка пустая. Значит, наши сетки различаются.} \\
            }
        }
    }
    & & \\ \\

    \theutterance \stepcounter{utterance}  
    & & & \multicolumn{2}{p{0.3\linewidth}}{
        \cellcolor[rgb]{0.9,0.9,0.9}{
            \makecell[{{p{\linewidth}}}]{
                \texttt{\tiny{[GM$|$GM]}}
                \texttt{continue} \\
            }
        }
    }
    & & \\ \\

    \theutterance \stepcounter{utterance}  
    & & \multicolumn{4}{p{0.6\linewidth}}{
        \cellcolor[rgb]{0.9,0.9,0.9}{
            \makecell[{{p{\linewidth}}}]{
                \texttt{\tiny{[P1$\langle$GM]}}
                \texttt{Теперь задайте вопрос, чтобы выяснить новые аспекты моей сетки, которые могут отличаться от вашей сетки. Начните с "ВОПРОС:" и не добавляйте ничего больше.} \\
            }
        }
    }
    & \\ \\

    \theutterance \stepcounter{utterance}  
    & \multicolumn{4}{p{0.6\linewidth}}{
        \cellcolor[rgb]{0.9,0.9,0.9}{
            \makecell[{{p{\linewidth}}}]{
                \texttt{\tiny{[P1$\rangle$GM]}}
                \texttt{ВОПРОС: В вашей верхней строке (строка 1) вторая клетка пустая — а какие значения в остальных клетках этой строки: столбец 1, столбец 3, столбец 4 и столбец 5?} \\
            }
        }
    }
    & & \\ \\

    \theutterance \stepcounter{utterance}  
    & & & \multicolumn{2}{p{0.3\linewidth}}{
        \cellcolor[rgb]{0.9,0.9,0.9}{
            \makecell[{{p{\linewidth}}}]{
                \texttt{\tiny{[GM$|$GM]}}
                \texttt{continue} \\
            }
        }
    }
    & & \\ \\

    \theutterance \stepcounter{utterance}  
    & & \multicolumn{4}{p{0.6\linewidth}}{
        \cellcolor[rgb]{0.9,0.9,0.9}{
            \makecell[{{p{\linewidth}}}]{
                \texttt{\tiny{[GM$\rangle$P2]}}
                \texttt{ОТВЕТ: Во второй строке моей сетки вторая клетка (столбец 2) пустая. Однако замечу различие: в моей верхней строке все клетки — X (X X X X X), тогда как у вас в верхней строке вторая клетка пустая. Значит, наши сетки различаются.} \\
\texttt{ВОПРОС: В вашей верхней строке (строка 1) вторая клетка пустая — а какие значения в остальных клетках этой строки: столбец 1, столбец 3, столбец 4 и столбец 5?Начните свой ответ с "ОТВЕТ:" и не добавляйте ничего другого.} \\
            }
        }
    }
    & \\ \\

    \theutterance \stepcounter{utterance}  
    & & & \multicolumn{4}{p{0.6\linewidth}}{
        \cellcolor[rgb]{0.9,0.9,0.9}{
            \makecell[{{p{\linewidth}}}]{
                \texttt{\tiny{[GM$\langle$P2]}}
                \texttt{ОТВЕТ: В моей верхней строке столбец 1 — X, столбец 3 — X, столбец 4 — X, столбец 5 — X, а столбец 2 пустой. Это отличается от вашей сетки, где вся верхняя строка заполнена X. Значит, наши сетки различаются.} \\
            }
        }
    }
     \\ \\

    \theutterance \stepcounter{utterance}  
    & & & \multicolumn{2}{p{0.3\linewidth}}{
        \cellcolor[rgb]{0.9,0.9,0.9}{
            \makecell[{{p{\linewidth}}}]{
                \texttt{\tiny{[GM$|$GM]}}
                \texttt{continue} \\
            }
        }
    }
    & & \\ \\

    \theutterance \stepcounter{utterance}  
    & & \multicolumn{4}{p{0.6\linewidth}}{
        \cellcolor[rgb]{0.9,0.9,0.9}{
            \makecell[{{p{\linewidth}}}]{
                \texttt{\tiny{[GM$\rangle$P2]}}
                \texttt{Теперь задайте вопрос, чтобы выяснить новые аспекты моей сетки, которые могут отличаться от вашей сетки. Начните с "ВОПРОС:" и не добавляйте ничего больше.} \\
            }
        }
    }
    & \\ \\

    \theutterance \stepcounter{utterance}  
    & & & \multicolumn{4}{p{0.6\linewidth}}{
        \cellcolor[rgb]{0.9,0.9,0.9}{
            \makecell[{{p{\linewidth}}}]{
                \texttt{\tiny{[GM$\langle$P2]}}
                \texttt{ВОПРОС: В вашей нижней строке (строка 5) все клетки заполнены X, или там тоже есть пустая клетка во втором столбце?} \\
            }
        }
    }
     \\ \\

    \theutterance \stepcounter{utterance}  
    & & & \multicolumn{2}{p{0.3\linewidth}}{
        \cellcolor[rgb]{0.9,0.9,0.9}{
            \makecell[{{p{\linewidth}}}]{
                \texttt{\tiny{[GM$|$GM]}}
                \texttt{continue} \\
            }
        }
    }
    & & \\ \\

    \theutterance \stepcounter{utterance}  
    & & \multicolumn{4}{p{0.6\linewidth}}{
        \cellcolor[rgb]{0.9,0.9,0.9}{
            \makecell[{{p{\linewidth}}}]{
                \texttt{\tiny{[P1$\langle$GM]}}
                \texttt{ОТВЕТ: В моей верхней строке столбец 1 — X, столбец 3 — X, столбец 4 — X, столбец 5 — X, а столбец 2 пустой. Это отличается от вашей сетки, где вся верхняя строка заполнена X. Значит, наши сетки различаются.} \\
\texttt{ВОПРОС: В вашей нижней строке (строка 5) все клетки заполнены X, или там тоже есть пустая клетка во втором столбце?Начните свой ответ с "ОТВЕТ:" и не добавляйте ничего другого.} \\
            }
        }
    }
    & \\ \\

    \theutterance \stepcounter{utterance}  
    & \multicolumn{4}{p{0.6\linewidth}}{
        \cellcolor[rgb]{0.9,0.9,0.9}{
            \makecell[{{p{\linewidth}}}]{
                \texttt{\tiny{[P1$\rangle$GM]}}
                \texttt{ОТВЕТ: В моей нижней строке (строка 5) все клетки заполнены X: X X X X X. У вас же во второй клетке нижней строки пусто. Значит, наши сетки различаются.} \\
            }
        }
    }
    & & \\ \\

    \theutterance \stepcounter{utterance}  
    & & & \multicolumn{2}{p{0.3\linewidth}}{
        \cellcolor[rgb]{0.9,0.9,0.9}{
            \makecell[{{p{\linewidth}}}]{
                \texttt{\tiny{[GM$|$GM]}}
                \texttt{continue} \\
            }
        }
    }
    & & \\ \\

    \theutterance \stepcounter{utterance}  
    & & \multicolumn{4}{p{0.6\linewidth}}{
        \cellcolor[rgb]{0.9,0.9,0.9}{
            \makecell[{{p{\linewidth}}}]{
                \texttt{\tiny{[P1$\langle$GM]}}
                \texttt{Теперь задайте вопрос, чтобы выяснить новые аспекты моей сетки, которые могут отличаться от вашей сетки. Начните с "ВОПРОС:" и не добавляйте ничего больше.} \\
            }
        }
    }
    & \\ \\

    \theutterance \stepcounter{utterance}  
    & \multicolumn{4}{p{0.6\linewidth}}{
        \cellcolor[rgb]{0.9,0.9,0.9}{
            \makecell[{{p{\linewidth}}}]{
                \texttt{\tiny{[P1$\rangle$GM]}}
                \texttt{ВОПРОС: Мы уже установили, что сетки различаются, но для полноты: в вашей сетке в третьей строке (строка 3) какие значения в столбцах 1, 2, 3, 4 и 5?} \\
            }
        }
    }
    & & \\ \\

    \theutterance \stepcounter{utterance}  
    & & & \multicolumn{2}{p{0.3\linewidth}}{
        \cellcolor[rgb]{0.9,0.9,0.9}{
            \makecell[{{p{\linewidth}}}]{
                \texttt{\tiny{[GM$|$GM]}}
                \texttt{continue} \\
            }
        }
    }
    & & \\ \\

    \theutterance \stepcounter{utterance}  
    & & \multicolumn{4}{p{0.6\linewidth}}{
        \cellcolor[rgb]{0.9,0.9,0.9}{
            \makecell[{{p{\linewidth}}}]{
                \texttt{\tiny{[GM$\rangle$P2]}}
                \texttt{ОТВЕТ: В моей нижней строке (строка 5) все клетки заполнены X: X X X X X. У вас же во второй клетке нижней строки пусто. Значит, наши сетки различаются.} \\
\texttt{ВОПРОС: Мы уже установили, что сетки различаются, но для полноты: в вашей сетке в третьей строке (строка 3) какие значения в столбцах 1, 2, 3, 4 и 5?Начните свой ответ с "ОТВЕТ:" и не добавляйте ничего другого.} \\
            }
        }
    }
    & \\ \\

    \theutterance \stepcounter{utterance}  
    & & & \multicolumn{4}{p{0.6\linewidth}}{
        \cellcolor[rgb]{0.9,0.9,0.9}{
            \makecell[{{p{\linewidth}}}]{
                \texttt{\tiny{[GM$\langle$P2]}}
                \texttt{ОТВЕТ: В моей третьей строке столбец 1 — X, столбец 2 — пусто, столбец 3 — пусто, столбец 4 — пусто, столбец 5 — X. Мы уже установили, что наши сетки различаются, так как в моей сетке второй столбец полностью пустой, а у вас внешняя рамка полностью заполнена X.} \\
            }
        }
    }
     \\ \\

    \theutterance \stepcounter{utterance}  
    & & & \multicolumn{2}{p{0.3\linewidth}}{
        \cellcolor[rgb]{0.9,0.9,0.9}{
            \makecell[{{p{\linewidth}}}]{
                \texttt{\tiny{[GM$|$GM]}}
                \texttt{continue} \\
            }
        }
    }
    & & \\ \\

    \theutterance \stepcounter{utterance}  
    & & \multicolumn{4}{p{0.6\linewidth}}{
        \cellcolor[rgb]{0.9,0.9,0.9}{
            \makecell[{{p{\linewidth}}}]{
                \texttt{\tiny{[GM$\rangle$P2]}}
                \texttt{Теперь задайте вопрос, чтобы выяснить новые аспекты моей сетки, которые могут отличаться от вашей сетки. Начните с "ВОПРОС:" и не добавляйте ничего больше.} \\
            }
        }
    }
    & \\ \\

    \theutterance \stepcounter{utterance}  
    & & & \multicolumn{4}{p{0.6\linewidth}}{
        \cellcolor[rgb]{0.9,0.9,0.9}{
            \makecell[{{p{\linewidth}}}]{
                \texttt{\tiny{[GM$\langle$P2]}}
                \texttt{ВОПРОС: Мы уже подтвердили различие, но для полноты: в вашей сетке во второй строке (строка 2) какие значения в столбцах 3, 4 и 5?} \\
            }
        }
    }
     \\ \\

    \theutterance \stepcounter{utterance}  
    & & & \multicolumn{2}{p{0.3\linewidth}}{
        \cellcolor[rgb]{0.9,0.9,0.9}{
            \makecell[{{p{\linewidth}}}]{
                \texttt{\tiny{[GM$|$GM]}}
                \texttt{continue} \\
            }
        }
    }
    & & \\ \\

    \theutterance \stepcounter{utterance}  
    & & \multicolumn{4}{p{0.6\linewidth}}{
        \cellcolor[rgb]{0.9,0.9,0.9}{
            \makecell[{{p{\linewidth}}}]{
                \texttt{\tiny{[P1$\langle$GM]}}
                \texttt{ОТВЕТ: В моей третьей строке столбец 1 — X, столбец 2 — пусто, столбец 3 — пусто, столбец 4 — пусто, столбец 5 — X. Мы уже установили, что наши сетки различаются, так как в моей сетке второй столбец полностью пустой, а у вас внешняя рамка полностью заполнена X.} \\
\texttt{ВОПРОС: Мы уже подтвердили различие, но для полноты: в вашей сетке во второй строке (строка 2) какие значения в столбцах 3, 4 и 5?Начните свой ответ с "ОТВЕТ:" и не добавляйте ничего другого.} \\
            }
        }
    }
    & \\ \\

    \theutterance \stepcounter{utterance}  
    & \multicolumn{4}{p{0.6\linewidth}}{
        \cellcolor[rgb]{0.9,0.9,0.9}{
            \makecell[{{p{\linewidth}}}]{
                \texttt{\tiny{[P1$\rangle$GM]}}
                \texttt{ОТВЕТ: В моей второй строке столбец 3 — пусто, столбец 4 — пусто, столбец 5 — X. Мы уже подтвердили, что наши сетки различаются: в моей сетке вся внешняя рамка заполнена X, тогда как у вас второй столбец полностью пустой (пустые клетки в строках 1 и 5 второго столбца).} \\
            }
        }
    }
    & & \\ \\

    \theutterance \stepcounter{utterance}  
    & & & \multicolumn{2}{p{0.3\linewidth}}{
        \cellcolor[rgb]{0.9,0.9,0.9}{
            \makecell[{{p{\linewidth}}}]{
                \texttt{\tiny{[GM$|$GM]}}
                \texttt{continue} \\
            }
        }
    }
    & & \\ \\

    \theutterance \stepcounter{utterance}  
    & & \multicolumn{4}{p{0.6\linewidth}}{
        \cellcolor[rgb]{0.9,0.9,0.9}{
            \makecell[{{p{\linewidth}}}]{
                \texttt{\tiny{[P1$\langle$GM]}}
                \texttt{Теперь задайте вопрос, чтобы выяснить новые аспекты моей сетки, которые могут отличаться от вашей сетки. Начните с "ВОПРОС:" и не добавляйте ничего больше.} \\
            }
        }
    }
    & \\ \\

    \theutterance \stepcounter{utterance}  
    & \multicolumn{4}{p{0.6\linewidth}}{
        \cellcolor[rgb]{0.9,0.9,0.9}{
            \makecell[{{p{\linewidth}}}]{
                \texttt{\tiny{[P1$\rangle$GM]}}
                \texttt{ВОПРОС: Для полноты картины: в вашей сетке в четвёртой строке (строка 4) какие значения в столбцах 1, 2, 3, 4 и 5?} \\
            }
        }
    }
    & & \\ \\

    \theutterance \stepcounter{utterance}  
    & & & \multicolumn{2}{p{0.3\linewidth}}{
        \cellcolor[rgb]{0.9,0.9,0.9}{
            \makecell[{{p{\linewidth}}}]{
                \texttt{\tiny{[GM$|$GM]}}
                \texttt{continue} \\
            }
        }
    }
    & & \\ \\

    \theutterance \stepcounter{utterance}  
    & & \multicolumn{4}{p{0.6\linewidth}}{
        \cellcolor[rgb]{0.9,0.9,0.9}{
            \makecell[{{p{\linewidth}}}]{
                \texttt{\tiny{[GM$\rangle$P2]}}
                \texttt{ОТВЕТ: В моей второй строке столбец 3 — пусто, столбец 4 — пусто, столбец 5 — X. Мы уже подтвердили, что наши сетки различаются: в моей сетке вся внешняя рамка заполнена X, тогда как у вас второй столбец полностью пустой (пустые клетки в строках 1 и 5 второго столбца).} \\
\texttt{ВОПРОС: Для полноты картины: в вашей сетке в четвёртой строке (строка 4) какие значения в столбцах 1, 2, 3, 4 и 5?Начните свой ответ с "ОТВЕТ:" и не добавляйте ничего другого.} \\
            }
        }
    }
    & \\ \\

    \theutterance \stepcounter{utterance}  
    & & & \multicolumn{4}{p{0.6\linewidth}}{
        \cellcolor[rgb]{0.9,0.9,0.9}{
            \makecell[{{p{\linewidth}}}]{
                \texttt{\tiny{[GM$\langle$P2]}}
                \texttt{ОТВЕТ: В моей четвёртой строке столбец 1 — X, столбец 2 — пусто, столбец 3 — пусто, столбец 4 — пусто, столбец 5 — X. Мы уже подтвердили, что наши сетки различаются: в моей сетке второй столбец полностью пустой, а у вас внешняя рамка полностью заполнена X.} \\
            }
        }
    }
     \\ \\

    \theutterance \stepcounter{utterance}  
    & & & \multicolumn{2}{p{0.3\linewidth}}{
        \cellcolor[rgb]{0.9,0.9,0.9}{
            \makecell[{{p{\linewidth}}}]{
                \texttt{\tiny{[GM$|$GM]}}
                \texttt{continue} \\
            }
        }
    }
    & & \\ \\

    \theutterance \stepcounter{utterance}  
    & & \multicolumn{4}{p{0.6\linewidth}}{
        \cellcolor[rgb]{0.9,0.9,0.9}{
            \makecell[{{p{\linewidth}}}]{
                \texttt{\tiny{[GM$\rangle$P2]}}
                \texttt{Теперь примите решение. Как вы думаете: ваша сетка и другая описанная сетка — одинаковые? Напишите "РЕШЕНИЕ: та же сетка", если вы считаете, что это одна и та же сетка, или "РЕШЕНИЕ: разные сетки", если вы считаете, что это разные сетки. Ничего больше не добавляйте.} \\
            }
        }
    }
    & \\ \\

    \theutterance \stepcounter{utterance}  
    & & & \multicolumn{4}{p{0.6\linewidth}}{
        \cellcolor[rgb]{0.9,0.9,0.9}{
            \makecell[{{p{\linewidth}}}]{
                \texttt{\tiny{[GM$\langle$P2]}}
                \texttt{РЕШЕНИЕ: разные сетки} \\
            }
        }
    }
     \\ \\

    \theutterance \stepcounter{utterance}  
    & & & \multicolumn{2}{p{0.3\linewidth}}{
        \cellcolor[rgb]{0.9,0.9,0.9}{
            \makecell[{{p{\linewidth}}}]{
                \texttt{\tiny{[GM$|$GM]}}
                \texttt{continue} \\
            }
        }
    }
    & & \\ \\

    \theutterance \stepcounter{utterance}  
    & & & \multicolumn{2}{p{0.3\linewidth}}{
        \cellcolor[rgb]{0.9,0.9,0.9}{
            \makecell[{{p{\linewidth}}}]{
                \texttt{\tiny{[GM$|$GM]}}
                \texttt{success} \\
            }
        }
    }
    & & \\ \\

    \theutterance \stepcounter{utterance}  
    & & \multicolumn{4}{p{0.6\linewidth}}{
        \cellcolor[rgb]{0.9,0.9,0.9}{
            \makecell[{{p{\linewidth}}}]{
                \texttt{\tiny{[P1$\langle$GM]}}
                \texttt{ОТВЕТ: В моей четвёртой строке столбец 1 — X, столбец 2 — пусто, столбец 3 — пусто, столбец 4 — пусто, столбец 5 — X. Мы уже подтвердили, что наши сетки различаются: в моей сетке второй столбец полностью пустой, а у вас внешняя рамка полностью заполнена X.} \\
\texttt{Теперь примите решение. Как вы думаете: ваша сетка и другая описанная сетка — одинаковые? Напишите "РЕШЕНИЕ: та же сетка", если вы считаете, что это одна и та же сетка, или "РЕШЕНИЕ: разные сетки", если вы считаете, что это разные сетки. Ничего больше не добавляйте.} \\
            }
        }
    }
    & \\ \\

    \theutterance \stepcounter{utterance}  
    & \multicolumn{4}{p{0.6\linewidth}}{
        \cellcolor[rgb]{0.9,0.9,0.9}{
            \makecell[{{p{\linewidth}}}]{
                \texttt{\tiny{[P1$\rangle$GM]}}
                \texttt{РЕШЕНИЕ: разные сетки} \\
            }
        }
    }
    & & \\ \\

    \theutterance \stepcounter{utterance}  
    & & & \multicolumn{2}{p{0.3\linewidth}}{
        \cellcolor[rgb]{0.9,0.9,0.9}{
            \makecell[{{p{\linewidth}}}]{
                \texttt{\tiny{[GM$|$GM]}}
                \texttt{continue} \\
            }
        }
    }
    & & \\ \\

    \theutterance \stepcounter{utterance}  
    & & & \multicolumn{2}{p{0.3\linewidth}}{
        \cellcolor[rgb]{0.9,0.9,0.9}{
            \makecell[{{p{\linewidth}}}]{
                \texttt{\tiny{[GM$|$GM]}}
                \texttt{success} \\
            }
        }
    }
    & & \\ \\

\end{supertabular}
}

\end{document}
